% !TEX root = ../main.tex

% 中英标题：\chapter{中文标题}[英文标题]
\chapter{预期目标}
\section{预期目标}
本课题的预期目标是基于多模态预训练的方法，成功实现蛋白质功能的准确预测，特别是在面对稀有功能标签和新发现的蛋白质时。通过利用蛋白质序列和结构特征，结合扩散模型与标签类条件分类器，我们期望能够提升蛋白质功能预测的精度和鲁棒性。主要的成果和目标可分为以下几点：

（1）针对蛋白质序列和结构数据的特点进行数据预处理，利用双分支变分自编码器（VAE）预训练并学习出适用于下游蛋白质功能预测的有效潜在空间表示。

（2）构建扩散模型，包括噪声添加和去噪两个过程，首先通过逐步添加高斯噪声模拟蛋白质数据的扩散，然后利用神经网络根据特定条件功能标签从噪声中恢复原始蛋白质表示。

（3）设计并实现基于扩散模型的标签类条件分类器，能够在无标签或稀有标签的数据条件下，精确预测蛋白质的功能。

（4）优化分类器使得其能够有效进行多标签分类，并在各种复杂条件下提供稳定的蛋白质功能预测结果，特别是针对序列同源相似性低的蛋白质。

（5）将上述研究工作整理成论文，争取以第一作者身份在JCR一区期刊上发表一篇具有影响力的学术论文。